
\section{Research questions and experimental setup}
\label{sec:simulation}
Given a diverse set of social scenarios, goals, and characters, we simulate agents' interactions. This is the first time that we could evaluate general, goal-oriented social agents in an interactive and systematic manner.
In the next three sections, we will demonstrate how \sotopia can be used to study these questions: (A) To which extent can we use GPT-4 \citep{openai2023gpt4} as a proxy for human judgment when it comes to evaluating agents' social interactions (\S\ref{sec:human_eval})? (B) What are the differences among models (\S\ref{sec:model_performance}) and between models and humans (\S\ref{sec:human_performance}) in their goal-oriented social intelligence? 

To study these questions, we create 40 agents, 90 relationships, and 90 scenarios following the generation procedure in \S\ref{sec:sotopia_characters_relationships}.
For each scenario, we sample 5 pairs of characters based on the scenario constraints, resulting in a set of 450 tasks.
For each task, we simulate the interaction between models by enumerating all model pairs. 
We also simulate the interaction between GPT-4 \citep{openai2023gpt4}\footnote{as will be shown in \S\ref{sec:model_performance} it is the best among models.} and humans on a challenging subset \sotopia-hard (\S\ref{sec:human_performance}) due to the limitation of resources. 

Specifically, we consider the following models for comparison: GPT-3.5 \citep{ouyang2022training}, GPT-4 \citep{openai2023gpt4}, Llama-2-70b-chat \citep{touvron2023llama}, and MPT-30b-chat \citep{MosaicML2023}. 
We set the temperature of the agents to 1 to encourage diversity of responses, and the temperature of the evaluator to 0 to ensure the stability of the evaluation. We use a fixed version of the above models to help reproducibility.\footnote{We fix GPT-4 to be \texttt{gpt-4-0613}, and GPT-3.5 to be \texttt{gpt-3.5-turbo-16k-0613}}
To use these models as agents in \sotopia, at each turn, we prompt the language model with the scenario, the character to play, and the interaction history to generate an action (see \S\ref{sec:episodes} for the possible actions). In this paper, as we are focusing on the use of \sotopia to understand social interaction, we use the prompt method for LLMs which is similar to the content of the interface for humans (Figure \ref{fig:human_performance_interface}). We leave leveraging novel prompting methods, e.g. Chain-of-Thought \citep{wei2022chain}, ReAct \citep{yao2022react}, as future work. 









